\chapter{System Design}

VisuaLearn AI is designed as a layered and modular system. The design separates user interaction, request orchestration, intelligent decision-making, educational content generation, simulation, visualization, voice interaction, and data persistence. This chapter presents the architecture, module-wise design, data flow, database design, and adaptive learning workflow.

\section[System Architecture]{\textbf{System Architecture}}
The system follows a five-layer architecture: presentation layer, application layer, intelligence layer, service layer, and data layer. Each layer has a clear responsibility and communicates with adjacent layers.

\begin{figure}[H]
\centering
\resizebox{0.85\textwidth}{!}{%
\begin{tikzpicture}[
box/.style={draw, rounded corners, minimum width=7.4cm, minimum height=0.9cm, align=center},
arrow/.style={-{Latex[length=2mm]}, thick},
node distance=0.6cm
]
\node[box] (p) {Presentation Layer: Chat, Learning Workspace, Voice, Settings};
\node[box, below=of p] (a) {Application Layer: Routing, Authentication, Sessions, Validation};
\node[box, below=of a] (i) {Intelligence Layer: Context Builder, Strategy Selector, Personalization};
\node[box, below=of i] (s) {Service Layer: Learning Content, Simulation, Visualization, Voice};
\node[box, below=of s] (d) {Data Layer: Users, Conversations, Resources, Feedback};
\draw[arrow] (p) -- (a);
\draw[arrow] (a) -- (i);
\draw[arrow] (i) -- (s);
\draw[arrow] (s) -- (d);
\draw[arrow] (d.east) .. controls +(2,0) and +(2,0) .. (a.east);
\end{tikzpicture}
}
\caption{Layered Architecture of VisuaLearn AI}
\label{fig:layeredarchitecture}
\end{figure}

\section[Presentation Layer Design]{\textbf{Presentation Layer Design}}
The presentation layer is implemented as a responsive web application. It provides the learner with the following views:
\begin{itemize}[leftmargin=1.2cm]
	\item Chat interface for asking questions and receiving responses.
	\item Learning workspace for notes, examples, flashcards, quizzes, mind maps, and simulations.
	\item Simulation viewer with playback controls.
	\item Visual widget frame for dynamic learning artifacts.
	\item Voice interface for spoken learning sessions.
	\item Settings page for personalization and account configuration.
\end{itemize}

The learning workspace is designed to reduce clutter. Instead of showing all generated content in one long page, each artifact is placed in a suitable tab or panel. This supports focused reading and easier revision.

\section[Application Layer Design]{\textbf{Application Layer Design}}
The application layer receives frontend requests and routes them to the correct backend module. It performs authentication, request validation, trace logging, error handling, and response formatting. It also supports streaming output so that learners can see responses progressively.

The application layer contains routes for chat, learning content, simulations, visualizations, voice sessions, documents, personas, progress, and settings. This modular routing structure makes the project easier to maintain and test.

\section[Intelligence Layer Design]{\textbf{Intelligence Layer Design}}
The intelligence layer is responsible for adapting the teaching strategy. It builds a context vector from learner signals and selects one of the supported response modes. The context includes:
\begin{itemize}[leftmargin=1.2cm]
	\item Cognitive state
	\item Engagement level
	\item Topic status
	\item Performance trend
	\item Confidence
	\item Topic difficulty
	\item Previous failures
	\item Response time
\end{itemize}

The teaching strategies are visual learning, guided steps, quiz checking, concise explanation, Socratic questioning, and remediation. The selected strategy influences how the learning response is generated and displayed.

\section[Service Layer Design]{\textbf{Service Layer Design}}
The service layer contains the core educational modules. The learning content service generates structured notes, examples, quizzes, flashcards, and mind maps. The simulation service determines whether a topic is procedural and generates state transitions. The visualization service prepares visual or spatial content. The voice service supports real-time spoken interaction. The resource service stores and retrieves generated learning artifacts.

\section[Data Layer Design]{\textbf{Data Layer Design}}
The data layer stores all persistent learning information. The main data entities are:
\begin{itemize}[leftmargin=1.2cm]
	\item Users: authentication and profile details.
	\item Conversations: learning sessions created by users.
	\item Messages: learner queries and system responses.
	\item Learning resources: notes, quizzes, flashcards, mind maps, simulations, and visual references.
	\item Feedback records: quiz results, completion data, and learner ratings.
	\item Adaptive decisions: selected strategy, context, and reward information.
\end{itemize}

\section[Data Flow Diagram]{\textbf{Data Flow Diagram}}
\begin{figure}[H]
\centering
\resizebox{0.95\textwidth}{!}{%
\begin{tikzpicture}[
entity/.style={draw, rectangle, rounded corners, minimum width=2.8cm, minimum height=0.8cm, align=center},
process/.style={draw, ellipse, minimum width=3.0cm, minimum height=0.95cm, align=center},
store/.style={draw, cylinder, shape border rotate=90, aspect=0.28, minimum width=1.5cm, minimum height=1.1cm, align=center},
arrow/.style={-{Latex[length=2mm]}, thick},
node distance=0.9cm
]
\node[entity] (user) {Learner};
\node[process, right=of user] (request) {Request Processing};
\node[process, right=of request] (engine) {Learning Engine};
\node[entity, right=of engine] (output) {Learning Output};
\node[store, below=of request] (db) {Data Store};
\draw[arrow] (user) -- (request);
\draw[arrow] (request) -- (engine);
\draw[arrow] (engine) -- (output);
\draw[arrow] (request) -- (db);
\draw[arrow] (engine) -- (db);
\draw[arrow] (db) -- (request);
\end{tikzpicture}
}
\caption{Data Flow Diagram of VisuaLearn AI}
\label{fig:dataflow}
\end{figure}

\section[Adaptive Learning Workflow]{\textbf{Adaptive Learning Workflow}}
The adaptive workflow connects learner behavior with future teaching decisions. When a learner asks a question, the system builds context and selects a response strategy. After the learner interacts with the output, feedback is stored. The feedback affects future decisions.

\begin{figure}[H]
\centering
\resizebox{0.95\textwidth}{!}{%
\begin{tikzpicture}[
box/.style={draw, rounded corners, minimum width=2.6cm, minimum height=0.85cm, align=center},
arrow/.style={-{Latex[length=2mm]}, thick},
node distance=0.55cm
]
\node[box] (q) {Question};
\node[box, right=of q] (c) {Context};
\node[box, right=of c] (p) {Policy};
\node[box, right=of p] (r) {Response};
\node[box, right=of r] (f) {Feedback};
\draw[arrow] (q) -- (c);
\draw[arrow] (c) -- (p);
\draw[arrow] (p) -- (r);
\draw[arrow] (r) -- (f);
\draw[arrow] (f.south) .. controls +(0,-0.7) and +(0,-0.7) .. (c.south);
\end{tikzpicture}
}
\caption{Adaptive Learning Workflow}
\label{fig:adaptiveworkflow}
\end{figure}

\section[Module Interaction]{\textbf{Module Interaction}}
The modules interact in a sequential but flexible manner. A single learner query may produce a direct explanation, a learning resource, a simulation, a visual widget, or a voice response. The backend decides which module is suitable based on query intent and learner context. All generated resources are saved so that the frontend can reload them during future sessions.

\section[Design Summary]{\textbf{Design Summary}}
The design of VisuaLearn AI emphasizes modularity, adaptability, reliability, and usability. The layered architecture allows independent development of the frontend, backend, intelligent services, and data storage. The adaptive workflow supports personalization, while the learning workspace organizes outputs into meaningful educational artifacts.
